\chapter{Wörterverzeichnis}

Dieser Abschnitt bildet das Wörterverzeichnis. Die Quellen der Definitionen werden abschnittsweise gegeben.

\section{Physikalische Grundlagen}
Die Definitionen dieses Abschnittes stammen teilweise frei aus Wikipedia.

\subsection{Das SI-Einheitensystem}
\begin{enumerate}
    \item   SI-Einheitensystem: das verbreitetste System für physikalische Einheiten. Alle Einheiten werden auf sieben Grundeinheiten Zeit, Länge, Masse, elektrische Stromstärke, thermodynamische Temperatur, Stoffmenge und Lichtstärke zurückgeführt. Zu jeder Grundeinheit wird eine physikalische Konstante festgelegt.
\end{enumerate}

\section{Begriffe aus der Chemie}
Die Definitionen dieses Abschnittes stammen teilweise frei aus Wikipedia und aus den Lehrbüchern Gerthsen \cite{gerthsen} und Demtröder \cite{demtröder2}.

\subsection{Begriffe und Größen}
\begin{enumerate}
    \item   Stoff: ein Element, eine Verbindung oder ein Gemisch mit bestimmten chemischen und physikalischen Eigenschaften.
    \item   Teilchenzahl N: die absolute Anzahl der Teilchen in einem System 
    \item   Stoffmenge n: relative Teilchenzahl, normiert an einer Grundmenge, die durch die Avogradokonstante festgelegt ist.
    \item   Mol: die SI-Einheit der Stoffmenge. 
    \item   Avogradokonstante: die Anzahl an Teilchen in einem Mol. Ein Mol eines Stoffes enthält  $6.02214076 \cdot 10^{23} \frac{1}{mol}$ Teilchen.
    \item   Molare Masse M: die Masse eines Stoffes pro Stoffmenge
    \item   Ion: ein elektrisch geladenes Atom oder Molekül. 
    \item   Anion: ein negativ geladenes Ion.
    \item   Kation: ein positiv geladenes Ion.
    \item   Oxidationszahl: die hypothetische Ionenladungen der Atome in einer chemischen Verbindung.
    \item   Wertigkeit von Ionen: die Anzahl der Elementarladungen eines Ions.
    \item   Faraday-Konstante:  die elektrische Ladung eines Mols einfach geladener Ionen.
\end{enumerate}

\subsection{Stoffe}
\begin{enumerate}
    \item   Lösung: ein homogenes Gemisch aus mindestens zwei chemischen Stoffen.
    \item   Schmelzung: der flüssige Aggragatszustand eines Stoffes, erreicht durch das Schmelzen dieses Stoffes.
    \item   Elektrolyte: Stoffe, deren Lösungen oder Schmelzen elektrischen Strom durch eine Ionenbewegung leiten.
\end{enumerate}

\subsection{Chemische Reaktionen}
\begin{enumerate}
    \item   Dissoziation: der Vorgang der Teilung eines Stoffes in zwei oder mehrere Moleküle, Atome oder Ionen.
    \item   Redoxreaktion: eine chemische Reaktion, wobei ein Reaktionsparter Elektronen abgibt und seine Oxidationszahl erhöht, während der andere Reaktionsparter diese Elektronen aufnimmt und damit seine Oxidationszahl erniedrigt.
    \item   Elektrolyse: eine durch ein Elektrolyt fließenden Strom ausgelöste Redoxreaktion innerhalb des Elektrolyts, die sogar zu dessen Dissoziation führen kann.
\end{enumerate}

\subsection{Stoffe, aus denen Elektrolyte bestehen können}
\begin{enumerate}
    \item   Säuren: chemische Verbindungen, die in der Lage sind, ein oder auch mehrere ihrer gebundenen H-Atome als Proton (H+) an einen Reaktionspartner zu übertragen, der für jedes zu bindende Proton ein freies Elektronenpaar zur Verfügung stellen muss.
    \item   Laugen/alkalische Lösungen: wässrige Lösungen, in denen die Konzentration der Hydroxid-Ionen OH- die der Oxonium-Ionen H30+ übersteigt, was einen pH-Wert der Lösung von größer als 7 zur Folge hat.
    \item   Basen: Verbindungen, die in wässriger Lösung in der Lage sind, Hydroxidionen (OH-) zu bilden und somit den pH-Wert einer Lösung zu erhöhen.
    \item   Salze: chemische Verbindungen, die aus elektrisch positiv geladenen Kationen und negativ geladenen Anionen aufgebaut sind.
    \item   Hydroxide:  salzähnliche Stoffe, die Hydroxid-Ionen (OH-) als negative Gitterbausteine (Anionen) enthalten.
\end{enumerate}